% 为什么只看树：三点网络 s->a (3,1)，a->t (2,1)，s->t (4,4)。
% (a) 流 (x_sa,x_at,x_st)=(1,1,2)，费用 1+1+8=10，三条弧都严格在界内，构成环。
%     沿环 s->a->t 增、s->t 减：每单位费用变化 1+1-4=-2，可推 min{3-1,2-1,2}=1。
% (b) 推送1单位后 (2,2,1)，费用8；a->t 到达上界，界内弧 s->a、s->t 不再成环。
\begin{tikzpicture}[>=Stealth,line width=.8pt,font=\fontsize{8}{10}\selectfont,
  v/.style={circle,draw=FigureStroke,fill=FigureBlueFill,line width=.8pt,minimum size=6.5mm,inner sep=0pt,text=FigureInk},
  free/.style={->,draw=FigureBlue,line width=1.5pt,shorten >=.5pt},
  bound/.style={->,draw=FigureStroke,line width=.8pt,dotted,shorten >=.5pt},
  lab/.style={text=FigureInk,inner sep=1.2pt,font=\fontsize{7.5}{9}\selectfont},
  inc/.style={text=FigureGreen,font=\fontsize{7.5}{9}\selectfont},
  note/.style={text=FigureMuted,font=\FigureNoteFont\fontsize{7.5}{9}\selectfont,align=center}]
% Panel labels share one baseline, 3 mm below the full row of content.
\begin{scope}[local bounding box=panel-content]
\begin{scope}
  \node[v] (s) at (0,0) {$s$}; \node[v] (a) at (1.9,1.45) {$a$}; \node[v] (t) at (3.8,0) {$t$};
  \draw[free] (s)--node[lab,above left=-1pt]{$1/3$\,\textcolor{FigureGreen}{$+$}} (a);
  \draw[free] (a)--node[lab,above right=-1pt]{$1/2$\,\textcolor{FigureGreen}{$+$}} (t);
  \draw[free] (s)--node[lab,below]{$2/4$\,\textcolor{FigureGreen}{$-$}} (t);
  \node[inc,align=center] at (1.9,.45) {每单位费用\\$1+1-4=-2$};
  \coordinate[overlay] (panel-a-center) at (1.9,-1.1);
\end{scope}
\begin{scope}[xshift=6.3cm]
  \node[v] (s) at (0,0) {$s$}; \node[v] (a) at (1.9,1.45) {$a$}; \node[v] (t) at (3.8,0) {$t$};
  \draw[free] (s)--node[lab,above left=-1pt]{$2/3$} (a);
  \draw[bound] (a)--node[lab,above right=-1pt]{$2/2$ 满} (t);
  \draw[free] (s)--node[lab,below]{$1/4$} (t);
  \coordinate[overlay] (panel-b-center) at (1.9,-1.1);
\end{scope}
\end{scope}
\node[font=\figurefont\fontsize{8}{10}\selectfont,text=FigureInk,anchor=north,inner sep=0pt,outer sep=0pt] at ([yshift=-3mm]panel-a-center |- panel-content.south) {(a)};
\node[font=\figurefont\fontsize{8}{10}\selectfont,text=FigureInk,anchor=north,inner sep=0pt,outer sep=0pt] at ([yshift=-3mm]panel-b-center |- panel-content.south) {(b)};
\end{tikzpicture}
